\section{2025-01-07}

\startframedtipp
{\bf Redundant}

In eine Tabelle tauchen die selben Informationen oefter auf
\stopframedtipp

\startitemize[packed]
\item {\bf Löschannomalie} wenn Daten X gelöscht werden und dadurch auch Daten Y verloren gehen, gelöcht werden.
\item {\bf Änderungsannomalie} Wenn die Daten an einem Punkt geändert werden, aber woanders, in der selben Tabelle, noch die Alten sind.
\item {\bf Einfügannomalie} Wenn die Datenbank um Attribute erweitert wird.
\stopitemize

Konsistenz - Zwei Dinge passen zusammen

\startframedtipp
{\bf Primaerschluessel}

Kann aus mehreren Attributen zusammengesetzt sein. Typischer weise aber soetwas wie {\em ID}. Er sollte sich allerdings an so wenig wie möglich Attributen  bedienen.
\stopframedtipp
